% !TEX root = ../../../main.tex
% Appendix entry: preprocessing.window_pool -- "Window pool" (Appendix A).
\entryfield{Algorithm}
\code{_run} in \code{Adapters/preprocessing_window_pool.py} maps the settings onto
\code{Working/training/pool.py}, the only implementation:
\begin{enumerate}
\item \emph{Reopen.} If \emph{pool} names a saved pool (id or name), \code{load_pool} reads it,
  checks its files against the key stored in the row, and the block returns it; nothing else
  runs.
\item \emph{The sets.} \emph{window\_sets} is a comma list of ids or names, or
  \texttt{all unlabelled} (every row \code{list_sets} reports as kind \texttt{unlabelled}).
  \code{set_windows} reads any kind of saved set as plain windows
  (recording, channel, start, length, rate): an unlabelled set, an earlier pool, a paired-job
  set or a one-channel set. A set that carries its own split loses it here.
\item \emph{The plan.} \code{plan_for} $\to$ \code{make_plan}: per recording
  \emph{identity} (\code{recording_identity} strips a \code{_fs<N>} suffix, so the fs1 and fs2
  files of one recording share one plan), the duration in seconds is cut by
  \code{blocked_split} into $B$ = n\_blocks equal blocks; the last
  $\mathrm{round}(B\cdot\text{test\_frac})$ (at least one when test\_frac $>0$) are test, the
  $\mathrm{round}(B\cdot\text{validation\_frac})$ before them validation, the rest train; a plan
  leaving no training block is refused. Consecutive blocks of one role merge into a
  \emph{stretch}, and gap\_min$\times 60$ seconds are carved off the start of every stretch
  that follows another role. Exam channels (\emph{exam\_channels}, one-based on the page, and
  the whole \emph{hold\_out\_pack}: A = CH1--4, B = CH5--8, C = CH9--12, D = CH13--16) are one
  exam stretch end to end. A plan that makes every channel an exam channel is refused.
\item \emph{Combine} (\code{combine}), in order:
  \begin{enumerate}
  \item refused (\code{PlanRefused}) if the longest window is longer than the gap;
  \item each window gets the role of the stretch it lies wholly inside
    (\code{assign_roles}); one in no stretch is dropped as \code{gap}, one across a stretch
    edge as \code{straddle};
  \item artifact spans are checked again (\code{_Exclusions}): any window touching a span a
    human marked \texttt{artifact} (annotations not soft-deleted, and window verdicts) or a
    span Settings excludes is dropped as \code{artifact};
  \item exact duplicates (same recording row, start, length) are dropped, the first set
    listed winning;
  \item within one scale on one recording identity and channel, overlapping windows are
    dropped, the first set listed winning (\code{overlap_within_scale});
  \item if per\_scale $> 0$: a seeded uniform sample of per\_scale windows per scale
    (\code{np.random.default_rng(seed)}), the rest \code{sampled_out};
  \item rule \texttt{no\_overlap} only: across scales, any window overlapping a kept window of a
    smaller scale is dropped (\code{overlap_across_scales}); under \texttt{within\_scale}
    overlap across scales is allowed, never across roles;
  \item \code{check_pool} recomputes every role from the plan and counts overlaps across
    roles on one recording's time (fs1 and fs2 together); any misplaced window or
    cross-role overlap raises \code{LeakageRefused}.
  \end{enumerate}
\item \emph{Save or reuse.} The pool's content key (\code{Pool.key}: a hash of the plan, rule,
  sample, seed, members and every window's file, channel, start, length and role) is looked
  up in \code{window_sets.recipe_hash}. A pool row with that key is reused; otherwise
  \code{save_pool} writes \code{pool_windows.npz} and \code{manifest.json} under
  \code{DATA/derived/window_sets/<name>_v<version>/} and a \code{window_sets} row with
  \code{window_pool_members} and \code{window_set_channels} rows.
\end{enumerate}

\entryfield{Assumptions}
\begin{itemize}
\item \emph{Units.} No signal values are read here; units do not enter.
\item \emph{Sampling rate.} Plans are in seconds and windows are compared on each recording's
  own time ($\text{start}/f_s$), so files at different rates combine correctly. The scale of
  a window is $\text{length}/f_s/60$ minutes, rounded to six decimals.
\item \emph{Missing values.} Non-finite windows are left out when an unlabelled set is built
  (\code{build_unlabelled_set}), not here. Windows from a labelled set were not screened
  for NaN \unverified{by this route}.
\item \emph{Packs.} hold\_out\_pack assumes the 16-channel, four-pack layout; a recording with
  fewer channels than the pack needs is refused.
\item \emph{The chain's span.} The block is a source (\code{source=True}): the validator
  refuses it anywhere but step 01, and the root signal is ignored. Because it reads the store
  (\code{conn}), it and every later step bypass the step cache.
\end{itemize}

\entryfield{Outputs}
A \texttt{WindowSet} (\code{pool_windowset} in \code{Working/training/shape.py}) whose starts
are the windows' own and whose feature table is metadata only: \code{pool_recording_id},
\code{pool_source_file}, \code{pool_channel}, \code{pool_start}, \code{pool_length},
\code{pool_fs}, \code{pool_scale_min}, \code{pool_role}, \code{pool_set_id}, the bounds of the
role stretch each window lies in (\code{pool_stretch_a}, \code{pool_stretch_b}) and the pool's
key (\code{pool_key}). Because a pool mixes lengths and rates, the WindowSet's own
\code{length} is the longest window and its \code{fs} the lowest rate; the true values are per
window. \code{meta["pool"]}: id, name, version, key, path, whether it was \code{saved},
\code{reused} or \code{re-opened}, counts by recording, scale and role, every drop by reason
(here and at set build), the checks, rule, sample, seed, members and the plan's stretches.

\entryfield{Cost}
The estimate is a fixed 10\,s: \code{combine} was measured at about 7\,s for six sets
(1.09 million windows) \unverified{on the development machine}. Memory is linear in the
number of windows offered. Not meant for the cluster.

\entryfield{Failure modes}
\begin{itemize}
\item \emph{No set ticked}, a set on a recording the plan does not cover, a gap shorter than
  the longest window, a pack on a recording with too few channels: refused with the reason.
\item \emph{A pool that leaks} across roles raises \code{LeakageRefused}; this should be
  unreachable, since roles come only from the plan.
\item \emph{Silent shrinkage.} A long gap\_min, many blocks or long windows drop many windows
  as \code{gap} or \code{straddle}. Nothing warns; read the drop counts.
\item \emph{A reused pool keeps its old name.} When the key matches an existing pool, the
  \emph{name} you typed is ignored.
\item \emph{Downstream blocks that do not read roles.} The roles ride as \code{pool_role}.
  \blockref{catalogue.shape_cluster} honours them; \blockref{catalogue.cluster} and
  \blockref{catalogue.shape_families} do not (see their entries).
\end{itemize}

\entryfield{Validation}
\code{check_pool} is a run-time guard, not a test: every saved pool has passed it. The tests
below check the plan, the drop reasons and the leakage refusals on small synthetic stores.

\entryfield{References}
The blocked-split rationale: \citet{roberts2017cross}. The Library side of window sets and
pools: \apppage{library.window-sets}. Research question 1:
\code{docs/rq_roundA/RQ1-cluster-vs-manual-labels.md}.

\entryfield{Tests}
\code{tests/test_training_pool.py}: sets ignore manual labels; artifact and Settings
exclusions; the held-out recording refused even with the unlock; hold-out pack; stretches
with a gap the same for every scale; a recording the plan does not cover refused; duplicates
and within-scale overlaps removed and counted; the strict rule (smaller scale wins); seeded
sampling; an artifact labelled after a set was saved still dropped; the fs1/fs2 pair never
on opposite sides; save and load round-trip. \code{tests/test_shape_chain.py}: the source may
only be step 01; the chain runs from saved sets to a kept tree; a saved pool reopens; no set
ticked is refused.
